\section{Introduction}
\label{sec:intro}

Automated program repair has moved from search- and template-based systems to
agents built on large language models that repair real projects by interacting
with an environment: localize a fault, propose an edit, run the test suite,
read the failure, revise~\cite{bouzenia2025repairagent,zhang2024autocoderover,%
yang2024sweagent,xia2025agentless,xia2024chatrepair}. That loop is
counterexample-guided repair by another name---each failing test is a
counterexample that ought to constrain the next proposal, exactly as a failing
input drives the synthesis loop of
\Sketch~\cite{solarlezama2006sketching}.

\paragraph{The cost of that loop is now measured.}
Analyzing $7{,}745$ SWE-bench agent traces and $3{,}000$ end-to-end repair
runs, Lin et al.~\cite{lin2026run} report that agents average $8.8$ test
executions per task, that the benefit concentrates on a minority of instances,
and that restricting execution entirely costs only $1.25$ percentage points of
resolve rate; their conclusion is that ``current agents apply execution
indiscriminately, paying their cost on instances where it provides little
benefit.'' Agent-authored tests likewise ``mimic a familiar
software-development practice while consuming interaction
budget''~\cite{chen2026rethinking}, and counting iteration honestly makes
self-repair's gains ``often modest''~\cite{olausson2024selfrepair}.

Three responses exist. Systems work makes each validation
cheaper~\cite{chen2021uniapr,xiao2024expressapr}; scaling work buys accuracy
with money, most explicitly CodeMonkeys at ${\approx}\$2{,}300$ for $57.4\%$ of
SWE-bench Verified~\cite{ehrlich2025codemonkeys}; and classical
generate-and-validate APR cuts the work itself, either by relating candidates
to each other~\cite{weimer2013leveraging,mechtaev2018testequivalence}, which
needs an enumerable candidate space an LLM does not supply, or by ordering the
tests so a doomed candidate dies on the
first~\cite{qi2013efficient,venugopal2020modification,lou2024when}, which needs
nothing of the candidate space at all. The second mechanism transfers. Our
question is what it should order from when the generator is a language model.

\paragraph{A refutation you already hold is such a criterion.}
If memory contains a counterexample $x$ that a candidate $p$ still fails, $p$
is refuted and validating it further buys nothing. In the synthesis loop this
idea comes from, the criterion could never fire---a constrained synthesizer
cannot emit a candidate that its accumulated counterexamples refute---but an
LLM proposer is under no such obligation and repeats itself in one proposal out
of four (\cref{sec:related-cegis}). The obstacle is that today's agents keep a
flat conversational transcript~\cite{xia2024chatrepair,shinn2023reflexion,%
chen2024selfdebug}, which neither \emph{retrieves}---it dilutes as the dialogue
grows---nor \emph{generalizes}: it records that patch $p_3$ failed test $t$,
not that every patch sharing $p_3$'s defect will fail the same way.

\paragraph{Our approach: keep the evidence, index it by where the patch edits.}
When a proposal $p$ is refuted by counterexample $x$, \CEGMem\ stores
$(p, x, \tau)$ where $\tau = \ttype(p,x)$ pairs $p$'s edit location with the
property the execution violated. The halves are available at different times,
and the mechanism rests on that: $\loc(p)$ is a diff against the buggy source
and needs no execution, while the violated property exists only after a failing
run. So the store is
\emph{indexed by $\loc$} and can be consulted for a candidate nobody has run.

That buys two operations. The agent can \emph{guard}: re-execute a stored
counterexample against the new candidate and, if it still fails, resolve the
round without calling the oracle---a decision that is verified rather than
predicted, since the candidate is blocked only after failing an input the
oracle itself produced. Inlined into
the oracle this is a prioritization discipline that stops on the first failing
case, and we claim no more than that; what is new is that the order comes from
counterexamples the loop discovered itself and kept across candidates. And it
can \emph{steer}: remove already-refuted classes from the proposer's support. Our study finds the first
carries essentially all of the value and the second none, which is not what we
set out to show.

\paragraph{This paper.}
We give \CEGMem\ a formal core (\cref{sec:theory}) and run it on real faults
rather than the simulation the design was first argued from: $99$ ConDefects
faults, a local \texttt{qwen2.5-coder:7b} proposer at $T = 1.0$, a differential
oracle over the shipped pool, five arms and four sweeps, \num{1485}
pre-specified main-grid cells and \num{3654} episodes. Two commitments make the
comparison tight: every arm is paired by \emph{common random
numbers}~\cite{law2000simulation}, so round~1 is the byte-identical draw in
every arm and later divergence is the treatment's own doing; and the primary
metric is pre-specified as repair rate at a matched \emph{oracle-call} budget,
not total calls, which an arm without a guard loses by construction.

Memory cuts oracle calls \fold{5.2} and, counting every guard replay as the
execution it is, total program executions \fold{2.2}, at unchanged repair rate;
at a matched two-call budget it lifts repair from $31.9\%$ to $55.8\%$. Against
a classical prioritizer replayed on identical candidates the execution saving
reverses and only the oracle-call reduction survives
(\cref{sec:results-policies}). Two
results cut against the architecture as proposed and we report them as primary:
prompt-side steering is null on every outcome metric while costing $+128\%$
prompt tokens and making the full agent the slowest of five arms, with no band
surviving a formal interaction test. The configuration we recommend is
consequently the guard alone.

\paragraph{Contributions.}
\emph{(i)} An architecture separating verification-side guarding from
prompt-side steering into independently analyzable operators
(\cref{sec:arch}).
\emph{(ii)} A formal core stated against a \emph{finite} oracle: guard
soundness, a $\Kreach{+}1$ round bound, strict redundancy reduction, and an
explicit account of what indexing does and does not buy (\cref{sec:theory}).
\emph{(iii)} A controlled execution of that loop on real faults under common
random numbers, with pre-specified metrics, falsifiers and a random-partition
control (\cref{sec:design,sec:results}).
\emph{(iv)} A negative result reported as a result: the prompt-side half does
not work at this scale, which relocates the value---and the
recommendation---onto the guard (\cref{sec:discussion}).
